\chapter{Weil Structures and Frobenius on Complete Complexes}
\label{chap:proetale-weil}

Let \(k=\Fq\), let \(\bar k=\Fbar\), and let
\(X_0\) be a separated scheme of finite type over \(k\).  Put
\[
  X=X_0\times_k\bar k.
\]
This chapter adds Frobenius descent data to the complete coefficient
categories.  No estimate on Frobenius eigenvalues is used.

\section{The two Frobenius descriptions}

Write
\[
 \sigma_q:a\longmapsto a^q
 \quad\text{and}\quad
 \varphi_q=\sigma_q^{-1}
\]
for arithmetic and geometric Frobenius in
\(\Gal(\bar k/k)\), respectively.  The element \(\varphi_q\) induces a
semilinear automorphism
\[
  \Fr_X:X\longrightarrow X.
\]
It is the geometric Frobenius automorphism used in every Weil structure
below.

There is also a scheme-theoretic description.  The absolute \(q\)-power
Frobenius
\[
  F^{(q)}_{X_0}:X_0\longrightarrow X_0
\]
is a universal homeomorphism.  Its pullback on the etale and pro-etale topoi
is therefore an equivalence by
\cref{thm:proetale-universal-homeomorphism}.  After geometric base change,
the contravariant pullback action agrees with the geometric-Frobenius
normalization.

\begin{convention}
The map on geometric points induced by raising coordinates to the
\(q\)-th power is called arithmetic Frobenius.  Cohomology is contravariant,
and throughout these notes the resulting operator on stalks and cohomology is
normalized as \emph{geometric} Frobenius \(\varphi_q=\sigma_q^{-1}\).
The invariant test is
\[
  \Fr\mid \Lambda(1)=q^{-1},
  \qquad
  \Fr\mid \Lambda(-1)=q.
\]
Any formula failing this test uses the opposite convention.
\end{convention}

\begin{proposition}[Frobenius pullback is invertible]
\label{prop:proetale-frobenius-equivalence}
Both \(\Fr_X^*\) and pullback by \(F^{(q)}_{X_0}\) are symmetric monoidal
autoequivalences of \(\Dcons\) and commute with derived completion, reduction
modulo \(\mathfrak m^n\), and scalar extension.
\end{proposition}

\begin{proof}
For \(\Fr_X\) this is immediate because it is an automorphism.  For
\(F^{(q)}_{X_0}\), use universal-homeomorphism invariance
\cite[Lemma~5.4.2]{BhattScholzeProEtale}.  The equivalence preserves the
coefficient sheaf and its finite quotients, so it commutes with the inverse
limit defining completion.
\end{proof}

\section{Weil complexes}

\begin{definition}[Integral Weil complex]
\label{def:proetale-weil-complex}
A constructible \(\Lambda\)-Weil complex on \(X_0\) is a pair
\[
  (K,\theta_K),\qquad
  K\in\Dcons(X,\Lambda),\qquad
  \theta_K:\Fr_X^*K\xrightarrow{\sim}K.
\]
A morphism \(u:(K,\theta_K)\to(L,\theta_L)\) satisfies
\[
  u\circ\theta_K
   =\theta_L\circ\Fr_X^*(u).
\]
The resulting category is denoted
\(\Dcons^W(X_0,\Lambda)\).  The same definition is used for
\(E\)- and \(\Qbarl\)-coefficients.
\end{definition}

Iteration of \(\theta_K\) gives compatible isomorphisms
\[
  (\Fr_X^n)^*K\xrightarrow{\sim}K
  \qquad(n\in\mathbf Z).
\]
Thus a Weil structure is a \(\mathbf Z\)-equivariant object.  It is essential
not to replace the isomorphism \(\theta_K\) by the assertion that
\(\Fr_X^*K=K\) definitionally.

\begin{proposition}[Descent and continuity]
\label{prop:weil-descent-continuity}
Pullback from \(X_0\) to \(X\) gives a fully faithful functor from
constructible complexes on \(X_0\) to Weil complexes on \(X_0\).  Its
essential image consists of those Weil complexes for which the
\(\mathbf Z\)-action generated by \(\theta_K\) extends continuously to
\[
  \Gal(\bar k/k)\simeq\widehat{\mathbf Z}.
\]
For an integral constructible complex, continuity can be checked after every
reduction modulo \(\mathfrak m^n\); on a finite stratification it is
equivalent to the finite-level Frobenius actions factoring through finite
quotients.
\end{proposition}

\begin{proof}
This is effective Galois descent on the pro-etale site.  A sheaf on \(X_0\)
pulls back with continuous descent data.  Conversely, continuous descent
data are effective because \(X\to X_0\) is a pro-etale Galois cover with
group \(\widehat{\mathbf Z}\).  By
\cref{thm:proetale-complete-system-equivalence}, an integral object is its
tower of finite reductions.  At each level the constructible Hom groups are
finite, so a continuous action factors through a finite quotient; compatible
finite-level descent reconstructs the complete object.
\end{proof}

\begin{warning}[Weil does not automatically mean continuous over \(E\)]
A single automorphism of an \(E\)-vector space defines a
\(\mathbf Z\)-action, but it need not extend to
\(\widehat{\mathbf Z}\).  For example, multiplication by a nonunit of
\(E\) preserves no \(\mathcal O_E\)-lattice as an automorphism.  Thus the
rational Weil category is larger than the category of objects descended
from \(X_0\).  The distinction disappears only after a Frobenius-stable
lattice or an equivalent continuity condition has been supplied.
\end{warning}

\section{Tate twists}

\begin{definition}[Integral Tate object]
\label{def:proetale-tate-object}
The pro-etale Tate sheaf is
\[
  \Zl(1):=\varprojlim_n \mu_{\ell^n}.
\]
For \(\Lambda=\mathcal O_E\), set
\[
 \Lambda(1)=\Zl(1)
   \widehat\otimes_{\Zl}\Lambda,
 \qquad
 \Lambda(r)=\Lambda(1)^{\widehat\otimes r}
\]
for \(r\ge0\), and define negative twists by the dual object.
For any \(K\), write \(K(r)=K\widehat\otimes\Lambda(r)\).
\end{definition}

\begin{proposition}[Frobenius normalization of twists]
\label{prop:tate-frobenius-normalization}
At a rational geometric point, geometric Frobenius acts on
\(\Lambda(r)\) by \(q^{-r}\).  Consequently twisting by \((r)\) lowers every
Frobenius weight by \(2r\).
\end{proposition}

\begin{proof}
Arithmetic Frobenius sends an \(\ell^n\)-th root of unity \(\zeta\) to
\(\zeta^q\).  Its inverse therefore acts through the inverse of \(q\) in
\((\mathbf Z/\ell^n)^{\times}\).  Passing to the inverse limit gives
\(q^{-1}\) on \(\Zl(1)\); tensor powers and duals give the general formula.
The assertion about weights is the definition of the twist normalization.
\end{proof}

\section{The six operations inherit Weil structures}

Let \(f_0:X_0\to Y_0\) be a separated finite-type morphism and
\(f:X\to Y\) its geometric base change.  The square
\[
\begin{tikzcd}
X \arrow[r,"\Fr_X"] \arrow[d,"f"'] & X \arrow[d,"f"] \\
Y \arrow[r,"\Fr_Y"'] & Y
\end{tikzcd}
\]
commutes.

\begin{theorem}[Weil-equivariant six operations]
\label{thm:proetale-weil-six-compatibility}
Every operation in
\cref{thm:proetale-six-functor-package} lifts canonically to Weil complexes.
For example, the Weil structure on \(\R f_!K\) is the composite
\[
 \Fr_Y^*\R f_!K
   \xrightarrow{\sim}
 \R f_!\Fr_X^*K
   \xrightarrow{\R f_!(\theta_K)}
 \R f_!K,
\]
and the analogous formula holds for \(\R f_*\).  Pullback, \(f^!\), completed
tensor, internal Hom, Verdier duality, Tate twists, base change, projection
formulas, and localization triangles are all Frobenius equivariant.
\end{theorem}

\begin{proof}
Apply the exchange equivalences of the six-functor formalism to the
commutative Frobenius square, then compose with the given Weil structure.
For tensor and internal Hom use the monoidality of \(\Fr^*\).  Coherence of
the six-functor package implies the cocycle and functoriality conditions.
All assertions may also be checked modulo \(\mathfrak m^n\), where they are
the standard torsion Weil formalism, and then recovered by completion.
\end{proof}

\begin{proetalebox}
The same complete object now carries topology, constructibility, all six
operations, and Frobenius.  There is no separate category of adic systems and
no later quotient by negligible towers.  This is the second major
pro-etale simplification.
\end{proetalebox}

\section{Changing coefficients without hiding lattices}

\begin{theorem}[Coefficient tower for Weil complexes]
\label{thm:proetale-weil-coefficient-tower}
Scalar extension gives functors
\[
 \Dcons^W(X_0,\Zl)
   \longrightarrow \Dcons^W(X_0,\mathcal O_E)
   \longrightarrow \Dcons^W(X_0,E)
   \longrightarrow \Dcons^W(X_0,\Qbarl).
\]
The essential image of
\[
 \Dcons^W(X_0,\mathcal O_E)[1/\ell]
   \longrightarrow \Dcons^W(X_0,E)
\]
consists exactly of the rational Weil complexes admitting a
Frobenius-stable \(\mathcal O_E\)-lattice.  Moreover
\[
 \Dcons^W(X_0,\Qbarl)
   =\varinjlim_{E/\Ql\text{ finite}} \Dcons^W(X_0,E).
\]
Every continuous rational Weil complex admits such a lattice after a finite
extension of \(E\), but the lattice is not canonical.
\end{theorem}

\begin{proof}
Scalar extension commutes with Frobenius pullback, so it transports the
structure isomorphism.  An integral Weil model plainly gives a stable
lattice.  Conversely a stable lattice, completed in the
\(\mathfrak m\)-adic topology, is an integral Weil object whose
localization is the given rational object.  A continuous representation has
compact image; a compact subgroup of \(\operatorname{GL}_n(E)\) stabilizes a
lattice after conjugation.  Apply this on finitely many strata and glue by
localization triangles.  The last formula follows from
\cref{eq:qbarl-colimit}, with the Frobenius isomorphism and each morphism
defined over a sufficiently large finite extension.
\end{proof}

\begin{formalization}
The coefficient API should return a proposition saying that a lattice
\emph{exists}; it should not choose one globally.  If a later construction
temporarily chooses a lattice, its output must be shown invariant under
replacing two lattices by a common finite-index sublattice.  The broader
rational Weil category and the stable-lattice subcategory should have
different names.
\end{formalization}

\section{Compactly supported cohomology and Frobenius}

Let \(a:X\to\Spec \bar k\) be the structural morphism.

\begin{theorem}[Adic perfectness]
\label{thm:proetale-weil-perfectness}
\label{thm:adic-perfectness}
For \(K\in\Dcons^W(X_0,\Lambda)\),
\begin{equation}
\label{eq:perfect-compact-cohomology}
 \R\Gamma_c(X,K):=\R a_!K
   \quad\text{belongs to}\quad \operatorname{Perf}(\Lambda).
\end{equation}
After scalar extension to \(E\), every
\(H_c^i(X,K_E)\) is finite-dimensional and vanishes outside a finite range.
The Weil structure induces an automorphism
\[
  F_K:\R\Gamma_c(X,K)\xrightarrow{\sim}\R\Gamma_c(X,K).
\]
If \(X_0\) is proper, \(\R\Gamma_c=\R\Gamma\).
\end{theorem}

\begin{proof}
By \cref{thm:proetale-weil-six-compatibility}, \(\R a_!K\) is a
constructible complete complex on the geometric point.  By
\cref{ex:constructible-point}, that category is \(\operatorname{Perf}(\Lambda)\).  The
Frobenius automorphism is the structure carried by \(\R a_!\); explicitly it
is
\[
 \R a_!K
   \xrightarrow{\sim}
 \R a_!\Fr_X^*K
   \xrightarrow{\R a_!(\theta_K)}
 \R a_!K.
\]
Properness identifies \(a_!\) and \(a_*\).
\end{proof}

\begin{formalization}
In the formalization \(\R a_!K\) is a Weil complex \(M\) on \(\Spec k_0\), and
\(H_c^r(X,K):=H^r(\R\Gamma(\Spec\bar k,M))\).  The Frobenius on it is transported along
the invariance \(\R\Gamma(\Spec\bar k,\sigma^*-)\cong\R\Gamma(\Spec\bar k,-)\) of derived
global sections under the geometric Galois Frobenius \(\sigma\) of \(\Spec\bar k\), which
follows from \(\sigma^*\cong\sigma^{-1}_*\) and the Leray isomorphism
\(\R\Gamma(\Spec\bar k,\R\sigma^{-1}_*-)\cong\R\Gamma(\Spec\bar k,-)\): it is
\[
 H^r(\R\Gamma(\Spec\bar k,M))
   \xrightarrow{\sim}
 H^r(\R\Gamma(\Spec\bar k,\sigma^*M))
   \xrightarrow{H^r(\varphi)}
 H^r(\R\Gamma(\Spec\bar k,M)),
\]
where \(\varphi:\sigma^*M\xrightarrow{\sim}M\) is the Weil structure of \(M\).
\end{formalization}

\begin{definition}[Cohomological trace and determinant]
For a rational Weil complex, define
\begin{align*}
 \Tr(F_K^n\mid\R\Gamma_c(X,K))
   &=\sum_i(-1)^i
     \Tr(F_K^n\mid H_c^i(X,K)),\\
 \det(1-TF_K\mid\R\Gamma_c(X,K))
   &=\prod_i
     \det(1-TF_K\mid H_c^i(X,K))^{(-1)^i}.
\end{align*}
The first expression is a scalar and the second is a normalized rational
function.  Both are well defined because of
\cref{thm:adic-perfectness}.
\end{definition}

\begin{corollary}[Compatibility of traces with coefficients]
\label{cor:weil-trace-coefficient-change}
The Frobenius traces, characteristic polynomials on individual cohomology
groups, and the alternating determinant commute with finite scalar extension
\(E\subset E'\) and hence define the corresponding invariants over
\(\Qbarl\).
\end{corollary}

\begin{proof}
Perfect base change identifies the cohomology after extension of scalars.
Trace and determinant of an endomorphism of a finite-dimensional vector
space are preserved by field extension.
\end{proof}

\begin{auditbox}
Perfectness is a finiteness statement, not a weight statement.  We have not
proved that Frobenius eigenvalues are algebraic, pure, integral, or
independent of \(\ell\), and we have not invoked the trace formula.  Those
assertions begin only after perverse sheaves and weights have been developed.
\end{auditbox}

\section{Why the next layer is perverse}

The six operations act on complexes, but the Weil proof needs an exact
category in which support dimension controls cohomological degree.
The middle-perversity t-structure supplies that category.  Once weights are
attached to its simple objects, the operations \(j_{!*}\), perverse
cohomology, vanishing cycles, and Fourier transform become exact or
one-sided-exact mechanisms that preserve sharp weight bounds.  The next
chapters exploit precisely this extra structure; ordinary cohomology sheaves
cannot provide the same local-to-global control.

\begin{powerbox}
The pro-etale site solves the coefficient problem; perversity solves the
weight-propagation problem.  Neither subsumes the other.  Their combination
is what lets the proof move from a local Frobenius estimate to global purity
without rebuilding adic limits at every step.
\end{powerbox}

